% CV of Xiaoyang Cao
% Build: cd cv && latexmk -xelatex CV_Xiaoyang_Cao.tex && cp CV_Xiaoyang_Cao.pdf ../files/CV_Xiaoyang_Cao.pdf && cp CV_Xiaoyang_Cao.pdf ../files/CV_Xiaoyang_Cao_2026.pdf
\documentclass[10pt,letterpaper]{article}

\usepackage{fontspec}            % Latin Modern under XeLaTeX
\usepackage[letterpaper,margin=0.78in,top=0.55in,bottom=0.55in]{geometry}
\usepackage{enumitem}
\usepackage{titlesec}
\usepackage[hidelinks]{hyperref}
\usepackage{xcolor}

\pagestyle{empty}
\setlength{\parindent}{0pt}
\setlength{\parskip}{0pt}
\linespread{1.0}
\raggedbottom

% ---------- section headings ----------
\titleformat{\section}{\large\bfseries\uppercase}{}{0pt}{}[\vspace{-6pt}\rule{\textwidth}{0.4pt}]
\titlespacing*{\section}{0pt}{9pt}{7pt}

% ---------- helpers ----------
\newcommand{\ulink}[2]{\href{#1}{\underline{#2}}}           % underlined link
\newcommand{\blink}[2]{[\href{#1}{\textbf{\underline{#2}}}]} % [arXiv]-style tag
\newcommand{\Bakker}{\ulink{https://miba.dev/}{Michiel Bakker}}
\newcommand{\Bayen}{\ulink{https://bayen.berkeley.edu/alex-bayen}{Alexandre Bayen}}
\newcommand{\sep}{\unskip\hspace{0.55em}$\cdot$\hspace{0.55em}\ignorespaces}
\newcommand{\me}{\textbf{Xiaoyang Cao}}

% Education entry: {institution}{location}{line(s) left}{dates}
\newcommand{\school}[4]{%
  \textbf{#1}\hfill #2\\
  \begin{minipage}[t]{0.68\textwidth}\itshape #3\end{minipage}%
  \hfill\begin{minipage}[t]{0.3\textwidth}\raggedleft #4\end{minipage}\par\vspace{5pt}}

% Research entry: {title + tags}{dates}{advisor line}
\newcommand{\project}[3]{%
  \par{\interlinepenalty=10000 % keep title and advisor lines together
  {#1}\hfill #2\\
  {\itshape #3}\par}\nopagebreak\vspace{1pt}}

\newenvironment{points}{%
  \begin{itemize}[leftmargin=1.1em,itemsep=1pt,parsep=0pt,topsep=2pt,label=\textbullet,
    before={\interlinepenalty=10000}]}% page breaks only between bullets
  {\end{itemize}\vspace{6pt}}

% Simple two-column row: left text, right-aligned date
\newcommand{\row}[2]{\par #1\hfill #2\par}

\begin{document}

% ================= header =================
\begin{center}
  {\fontsize{17.2}{20}\selectfont\bfseries Xiaoyang Cao}\\[4pt]
  Cambridge, MA\sep (+1) 617-216-8184\sep
  \ulink{mailto:xycao@mit.edu}{xycao@mit.edu}\sep
  \ulink{https://xiaoyangcao1113.github.io/}{Personal Website}\sep
  \ulink{https://scholar.google.com/citations?user=XVWuYHAAAAAJ}{Google Scholar}\sep
  \ulink{https://github.com/XiaoyangCao1113}{GitHub}
\end{center}
\vspace{-6pt}

% ================= education =================
\section{Education}

\school{Massachusetts Institute of Technology}{Cambridge, MA}
  {M.S. in Computational Science \& Engineering\\
   Concurrent M.S. in Technology \& Policy\\
   Laboratory for Information and Decision Systems (LIDS)\\
   Research advisor: Prof. \Bakker}
  {Sep 2025 -- May 2027 (expected)}

\school{Tsinghua University}{Beijing, China}
  {B.S. in Mathematics and Physics}
  {Sep 2021 -- Jun 2025}

\school{University of California, Berkeley}{Berkeley, CA}
  {Exchange Student}
  {Jan 2024 -- Aug 2024}

% ================= interests =================
\section{Research Interests}

How intelligence emerges in collectives of interacting agents (multi-agent reinforcement learning, populations of LLM agents), and how oversight can keep pace with increasingly capable AI systems (scalable oversight, alignment).

% ================= publications =================
\section{Publications}

\begin{enumerate}[label={[\arabic*]},leftmargin=1.8em,itemsep=4pt,parsep=0pt,topsep=0pt]
  \item \me, Jingqi Li, Zhe Fu, Alexandre M. Bayen. ``Who Bears the Burden? Learning Responsibility for Shared Constraints in Multi-Agent Reinforcement Learning.'' \textit{Under review}, 2026.
  \item \me, Siddarth Srinivasan, Michiel A. Bakker. ``Do Modules Stay in Their Lane? Role Drift in Compound LLM Systems.'' \textit{Under review}, 2026.
  \item \me, Zelai Xu, Mo Guang, Kaiwen Long, Michiel Bakker, Yu Wang, Chao Yu. ``RE-PO: Robust Enhanced Policy Optimization as a General Framework for LLM Alignment.'' \textit{International Conference on Learning Representations (ICLR)}, 2026.
  \item \me, Zhe Fu, Alexandre M. Bayen. ``Pareto Control Barrier Function for Inner Safe Set Maximization Under Input Constraints.'' \textit{American Control Conference (ACC)}, 2025.
  \item \me, Dingyi Zhuang, Shenhao Wang, Jinhua Zhao. ``Virtual Nodes Improve Long-term Traffic Prediction.'' \textit{Transportation Research Board Annual Meeting (TRB)}, 2025.
\end{enumerate}

% ================= research experience =================
\section{Research Experience}

\project{\textbf{Learning Responsibility for Shared Constraints in MARL} \textbf{[1]}
  \blink{https://arxiv.org/abs/2610.07491}{arXiv} \blink{https://lira-marl.github.io/}{Project}}
  {May 2025 -- Sep 2026}{Advised by Prof. \Bayen, UC Berkeley}
\begin{points}
  \item Proposed LiRA (Lagrangian Responsibility Allocation): a common multiplier enforces a shared cost budget, while learned per-agent responsibility shares redistribute its penalty by optimizing social welfare, without modifying rewards or constraints.
  \item Proved that for convex games varying the shares induces a smooth family of normalized generalized Nash equilibria in which active constraints stay at budget while welfare varies, and derived a welfare gradient that accounts for both learning updates and the induced shift in data distribution.
  \item Improved average social welfare by up to 29\% over uniform and agent-specific multiplier baselines across CityLearn, MABIM, Harvest, and MetaDrive (3 to 400 agents).
\end{points}

\project{\textbf{Role Drift in Compound LLM Systems} \textbf{[2]}
  \blink{https://arxiv.org/abs/2607.21627}{arXiv}
  \blink{https://venturebeat.com/ai/one-ai-module-faked-86-of-a-pipelines-accuracy-gains-by-feeding-another-the-answers}{VentureBeat}}
  {Apr 2026 -- Sep 2026}{Advised by Prof. \Bakker, MIT}
\begin{points}
  \item Identified \textit{role drift}: end-to-end reinforcement learning can improve terminal performance while causing specialized modules to abandon their intended functions.
  \item Developed role-level evaluations showing that 86\% of an apparent 31-percentage-point accuracy gain disappeared once the decomposer's assigned function was measured, demonstrating that end-to-end metrics can conceal internal failure.
  \item Proposed Role Anchor, a training objective that preserves the influence of role specifications while allowing modules to improve under end-to-end optimization.
\end{points}

\project{\textbf{RE-PO: Robust LLM Preference Alignment} (ICLR 2026) \textbf{[3]}
  \blink{https://arxiv.org/abs/2509.24159}{arXiv} \blink{https://repo-alignment.github.io/}{Project}}
  {Mar 2025 -- Sep 2025}{Advised by Prof. \ulink{https://web.ee.tsinghua.edu.cn/wangyu/en/}{Yu Wang}, Tsinghua University}
\begin{points}
  \item Proposed an EM-based framework that models preference-label correctness as a latent variable, jointly estimates label and annotator reliability, and systematically robustifies DPO, IPO, SimPO, and CPO.
  \item Proved convergence to the true noise level under calibration and improved AlpacaEval 2 win rates by up to 7.0 percentage points on Mistral-7B and Llama-3-8B, with additional gains on Arena-Hard and MultiPref.
\end{points}

\project{\textbf{Safe Control via Pareto Control Barrier Functions} (ACC 2025) \textbf{[4]}
  \blink{https://arxiv.org/abs/2410.04260}{arXiv}}
  {May 2024 -- Oct 2024}{Advised by Prof. \Bayen, UC Berkeley}
\begin{points}
  \item Developed the Pareto Control Barrier Function algorithm by integrating Pareto multi-task learning into neural control barrier functions, maximizing the inner safe set while guaranteeing safety under input constraints.
  \item Matched the Hamilton--Jacobi viability kernel on a 2D inverted pendulum and achieved a 700\% larger inner safe set than the benchmark on a 12-D quadrotor obstacle-avoidance task.
\end{points}

\project{\textbf{Virtual Nodes for Long-term Traffic Prediction} (TRB 2025) \textbf{[5]}
  \blink{https://arxiv.org/abs/2501.10048}{arXiv}}
  {May 2024 -- Sep 2024}{Advised by Prof. \ulink{https://zhaojinhua.com/}{Jinhua Zhao}, MIT}
\begin{points}
  \item Introduced virtual nodes into spatio-temporal graph neural networks to mitigate over-squashing in long-horizon traffic forecasting, using a semi-adaptive adjacency matrix learned from node embeddings.
  \item Reduced RMSE by 6.27\% and MAPE by 5.04\% on a five-year, 716-sensor San Diego dataset; the learned adjacency emphasized traffic-intensive intersections, improving interpretability.
\end{points}
\vspace{-6pt}

% ================= awards =================
\section{Awards \& Honors}

\row{\textbf{Outstanding Winner \& INFORMS Award}, Mathematical Contest in Modeling -- top 1\%}{May 2023}
\row{\textbf{Outstanding Innovation in Science \& Technology Scholarship}, Tsinghua University}{2023, 2024}
\row{\textbf{Comprehensive Excellence Scholarship}, Tsinghua University}{Oct 2023}
\row{\textbf{Outstanding Academic Scholarship}, Tsinghua University}{Oct 2022}

% ================= service =================
\section{Service \& Teaching}

\row{\textbf{Conference Reviewer:} ICLR 2027, CoRL 2026}{}
\row{\textbf{Teaching Assistant:} Calculus A, Tsinghua University}{Spring 2025}

% ================= skills =================
\section{Skills \& Languages}

\textbf{ML \& RL:} Python, PyTorch, Hugging Face Transformers / TRL, vLLM, DeepSpeed, FSDP, PettingZoo, RLlib\\
\textbf{Compute:} Multi-GPU, multi-node training on H100 / H200 / B200 clusters\\
\textbf{Languages:} English (fluent), Chinese (native)

\end{document}
